\section{Method}
\label{sec:method}

\subsection{Problem formulation}
Let $f_m$ be model $m$ with greedy decoding and let $x_i$ be the canonical prompt for item $i$.
A perturbation $p$ is a string edit that preserves meaning and is invisible or nearly invisible when rendered.
It produces $p(x_i)$.
We measure divergence relative to the model's own canonical answer, not relative to the gold label:
\begin{equation}
\label{eq:div}
D_m(p) = \frac{1}{|\gI_p|}\sum_{i \in \gI_p} \1\big[g(f_m(p(x_i))) \neq g(f_m(x_i))\big],
\end{equation}
where $g$ maps a raw output to the quantity of interest (a parsed label, the extracted code, or a unit-test outcome) and $\gI_p = \{i : p(x_i) \neq x_i\}$ is the set of items the edit actually changes.
The \ctrl condition sets $p$ to the identity but places the prompt at a different position in the batch, so $D_m(\ctrl)$ is the noise floor of inference.

\subsection{Perturbations}
\Tabref{tab:perturbations} lists the edits.
The whitespace family has four edits shared by all tasks, plus no-break spaces (\cond{nbsp}) for classification or tab indentation (\cond{tabs}) for code.
The punctuation family replaces ASCII quotes, apostrophes, dashes and ellipses with their typographic Unicode forms, separately and all together (\cond{punct_all}).
We wrote the task instructions to contain double quotes, an apostrophe (``Don't'') and a spaced hyphen, so that every punctuation edit changes every prompt.
The exception is \cond{ellipsis_char}, which applies only to the 40 classification and 9 code prompts that contain \texttt{...}.
Edits that leave a prompt unchanged are excluded from all rates.

Code prompts must stay valid Python.
We apply \cond{double_space} and all punctuation edits only to natural language: the instruction, docstrings and comments, which we locate with Python's \texttt{tokenize} module.
Doctest lines (\texttt{>>>} and their expected output) are left untouched.
Layout edits (\cond{trailing_space}, \cond{crlf}, \cond{tabs}, \cond{trailing_newlines}) apply to the whole message.

\begin{table}[t]
    \centering
    \small
    \resizebox{\textwidth}{!}{%
    \begin{tabular}{@{}llp{8.4cm}l@{}}
        \toprule
        \textbf{Family} & \textbf{Condition} & \textbf{Edit} & \textbf{Applies to} \\
        \midrule
        Control & \ctrl & Byte-identical to canonical, different batch position & all \\
        \midrule
        \multirow{6}{*}{Whitespace} & \cond{trailing_space} & Append a space to every line & all \\
        & \cond{crlf} & Replace \texttt{\textbackslash n} with \texttt{\textbackslash r\textbackslash n} & all \\
        & \cond{double_space} & Replace every single inter-word space with two spaces & all (code: prose only) \\
        & \cond{trailing_newlines} & Append \texttt{\textbackslash n\textbackslash n} to the message & all \\
        & \cond{nbsp} & Replace inter-word spaces with no-break space U+00A0 & classification \\
        & \cond{tabs} & Replace 4-space indentation with a tab & \humaneval \\
        \midrule
        \multirow{5}{*}{Punctuation} & \cond{curly_quotes} & \texttt{"x"} $\rightarrow$ \textquotedblleft x\textquotedblright & all (code: prose only) \\
        & \cond{curly_apostrophe} & \texttt{'} $\rightarrow$ \textquoteright\ / \textquoteleft & all (code: prose only) \\
        & \cond{unicode_dash} & \texttt{~-~} $\rightarrow$ en dash, \texttt{--} $\rightarrow$ em dash, digit ranges $\rightarrow$ en dash & all (code: prose only) \\
        & \cond{ellipsis_char} & \texttt{...} $\rightarrow$ U+2026 (\ldots) & 40 cls / 9 code prompts \\
        & \cond{punct_all} & All four punctuation edits combined & all (code: prose only) \\
        \bottomrule
    \end{tabular}
    }
    \caption{The ten invisible formatting edits plus the identical-prompt control. ``Prose only'' means the edit touches the instruction, docstrings and comments but never code tokens or doctest lines, so programs stay valid.}
    \label{tab:perturbations}
\end{table}


\subsection{Tasks and models}
\para{Classification.}
We sample 400 items each (seed 42) from the GLUE~\cite{wang2018glue} validation split of \sst~\cite{socher2013recursive}, the test split of \agnews~\cite{zhang2015character} and the validation split of \boolq~\cite{clark2019boolq} (passages of at most 180 words).
Each prompt is a zero-shot instruction such as \texttt{Classify the sentiment of the movie review below as "positive" or "negative".} followed by \texttt{Don't explain - answer with one word only.} and the input.
We parse the earliest label keyword in the output, map synonyms, and assign \texttt{other} to anything else.

\para{Code generation.}
We use all 164 \humaneval problems~\cite{chen2021evaluating} and the 300-problem \mbpp test split~\cite{austin2021program}.
\humaneval prompts ask the model to complete the given function and return it in one fenced block. \mbpp prompts give the task text and its three \texttt{assert} tests.
We extract the first fenced block that contains a \texttt{def} and normalize CRLF and trailing whitespace.
We then run the unit tests in an isolated subprocess (\texttt{python -I}, 2\,GB memory limit, 20\,s timeout).

\para{Models.}
We run 11 open instruction-tuned models: \qwentwo 0.5B, 1.5B, 3B, 7B and 14B~\cite{qwen2025qwen25}; \qwenthree 0.6B, 4B and 8B~\cite{yang2025qwen3}, with thinking mode disabled; and Llama-3.2 1B and 3B and Llama-3.1 8B~\cite{grattafiori2024llama3}.
Most checkpoints are AWQ 4-bit~\cite{lin2024awq}. Llama-3.2-1B, \qwenthree-0.6B and \qwenthree-8B run in bf16.
All models run in vLLM 0.32~\cite{kwon2023efficient} on NVIDIA RTX A6000 GPUs with temperature 0, using at most 12 new tokens for classification and 512 for code.
We call the four models under 3B parameters \emph{small} and the other seven \emph{large}.
Each model sees 19,968 prompts (1,664 items $\times$ 12 conditions), for 239,616 generations in total including a replication run.
We had no commercial API access, so the study covers open-weight models only.

\subsection{Metrics}
We report four divergence metrics, each an instance of \eqref{eq:div}:
(1) \textbf{label flip}, where the parsed classification label differs;
(2) \textbf{code changed}, where the normalized extracted code differs;
(3) \textbf{pass/fail flip}, where the unit-test outcome differs; and
(4) \textbf{raw text changed}, where the output string differs.
We also report accuracy against gold labels and tests.
For every model and prompt, we apply the chat template and tokenizer and record whether the token sequence changes and how many positions differ (\dtok, computed with \texttt{difflib}).

\subsection{Statistical analysis}
All metrics are paired at the item level.
We compute 95\% confidence intervals with a bootstrap over items (2,000 resamples). For the small/large ratio, the bootstrap resamples the same items for all models.
For each of the 110 model $\times$ condition cells, we test perturbation against control with an exact McNemar test on discordant items and apply Holm correction~\cite{holm1979simple}.
We fit GEE logistic models~\cite{liang1986longitudinal} with exchangeable correlation clustered by item, of the form flip $\sim$ punctuation + small + task, with and without $\log(1+\Delta\text{tok})$, using statsmodels~\cite{seabold2010statsmodels}.

\para{Reproducibility checks.}
We re-ran \qwentwo-0.5B in a separate engine process on a different GPU. All 19,968 outputs were byte-identical to the first run, so every rate for that model replicates exactly.
Llama-3.x chat templates strip trailing whitespace from messages, so \cond{trailing_newlines} leaves their token sequence unchanged. This gives a natural placebo: an edit to the raw string that the model never sees.
